\begin{table}[h]
    
\noindent
\begin{minipage}{0.15\textwidth}
\centering
\begin{tabular}{c | c}
    $T_1^1$ & $T_1^0$ \\
    $\wt(s_1,1)$~& ~$\wt(s_1,1)$ \\ 
    $\wt(s_1,2)$~& ~$\wt(s_1,2)$ \\ 
    $\wt(v_1, 1)$~& ~$\wt(v_1, 0)$ \\

\end{tabular}
\vspace{.1cm}

\begin{tabular}{c | c}
    $T_2^1$ & $T_2^0$ \\
    $\wt(s_2,1)$~& ~$\wt(s_2,1)$ \\ 
    $\wt(s_2,2)$~& ~$\wt(s_2,2)$ \\ 
    $\wt(v_2, 1)$~& ~$\wt(v_2, 0)$ \\
\end{tabular}
\vspace{.1cm}

\begin{tabular}{c | c}
    $T_3^1$ & $T_3^0$ \\
    $\wt(s_3,1)$~& ~$\wt(s_3,1)$ \\ 
    $\wt(s_3,2)$~& ~$\wt(s_3,2)$ \\ 
    $\wt(v_3, 1)$~& ~$\wt(v_3, 0)$ \\
   
\end{tabular}

\end{minipage}
\hfill
\begin{minipage}{0.15\textwidth}
\centering
\centering
\begin{tabular}{c | c}
    $T_{x_1}$ & $T_{\neg x_1}$ \\
    $\rd(v_1,1)$ ~& ~$\rd(v_1,0)$ \\ 
    $\wt(c_1,1)$~ &  \\ 
    $\wt(c_2, 1)$~&  \\

\end{tabular}
\vspace{.1cm}

\begin{tabular}{c | c}
   $T_{x_2}$ & $T_{\neg x_2}$ \\
    $\rd(v_2,1)$ ~& ~$\rd(v_2,0)$ \\ 
    $\wt(c_1,1)$~ & ~$\wt(c_2,1)$ \\ 
    
\end{tabular}
\vspace{.1cm}

\begin{tabular}{c | c}
    $T_{x_3}$ & $T_{\neg x_3}$ \\
    $\rd(v_3,1)$ ~& ~$\rd(v_3,0)$ \\ 
    $\wt(c_1,1)$~ & ~$\wt(c_2,1)$ \\ 
    
   
\end{tabular}
\end{minipage}
\hfill
\begin{minipage}{0.15\textwidth}
\centering
\begin{tabular}{c}
    $T_f$ \\
    $\rd(c_1, 1)$ \\
    $\rd(c_2, 1)$ \\
    $\rd(s_1, 2)$ \\
    $\rd(s_1, 1)$ \\
    $\rd(s_2, 2)$ \\
    $\rd(s_2, 1)$ \\
    $\rd(s_3, 2)$ \\
    $\rd(s_3, 1)$ \\

\end{tabular}
\end{minipage}
\vspace{2em}
\caption{$\varphi=(x_1\lor x_2\lor x_3)\land (x_1 \lor \neg x_2\lor \neg x_3)$}
    \label{tab:example-finegrainedwra}
\end{table}

Table~\ref{tab:example-finegrainedwra} gives an instance of the construction of $\expartial$ for the fine-grained hardness reduction for \threewriter~ and (Figure~\ref{example-Execution-ra}) gives the complete execution graph for it. The execution graph is constructed for $x_1=1, x_2=0, x_3=0$, a satisfying assignment for $\varphi$. The $\rf$ and $\mo$-edges follow the rule of construction for them given in the reduction we discussed before. The $\rf$ synthesis is as follows:
\begin{itemize}
    \item $\tuple{T_{x_1}:\wt(c_1,1)}~\rf~\tuple{T_f:\rd(c_1,1)}$, $\tuple{T_{x_1}:\wt(c_2,1)}~\rf~\tuple{T_f:\rd(c_2,1)}$
    \item $\tuple{T_1^1:\wt(s_1,2)}~\rf~\tuple{T_f:\rd(s_1,2)}$, $\tuple{T_1^0:\wt(s_1,1)}~\rf~\tuple{T_f:\rd(s_1,1)}$
    \item $\tuple{T_2^0:\wt(s_2,2)}~\rf~\tuple{T_f:\rd(s_2,2)}$, $\tuple{T_2^1:\wt(s_2,1)}~\rf~\tuple{T_f:\rd(s_2,1)}$
    \item $\tuple{T_3^0:\wt(s_3,2)}~\rf~\tuple{T_f:\rd(s_3,2)}$, $\tuple{T_3^1:\wt(s_3,1)}~\rf~\tuple{T_f:\rd(s_3,1)}$
\end{itemize}

All other $\rf$-edges are trivial. They are not reflected in Figure~\ref{example-Execution-ra}. Let us construct the $\mo$-edges.
\begin{itemize}
 \item $\tuple{T_1^1:\wt{s_1,2}} \xra{\mo}\allowbreak \tuple{T_1^0:\wt{s_1,1}}$ , $\tuple{T_1^1:\wt{s_1,1}} \xra{\mo}\allowbreak \tuple{T_1^0:\wt{s_1,1}}$
 
 \item $\tuple{T_j^0:\wt{s_j,2}} \xra{\mo}\allowbreak \tuple{T_j^1:\wt{s_j,1}}$, $\tuple{T_j^0:\wt{s_j,1}} \xra{\mo}\allowbreak \tuple{T_j^1:\wt{s_j,1}}$ for $j=2,3$
 
 \item $\tuple{T_{x_j}:\wt{c_1,1}} \xra{\mo}\allowbreak \tuple{T_{x_1}:\wt{c_1,1}}$, $\tuple{T_{x_j}:\wt{c_2,1}} \xra{\mo}\allowbreak \tuple{T_{x_1}:\wt{c_2,1}}$ for $j=1,2$.
\end{itemize}

Other $\mo$ edges follow the direction of $\po$ and are not explicitly constructed in Figure~\ref{example-Execution-ra}. %Note that the resulting execution graph satisfies all the memory axioms of $\sramm$ (Refer Table~\ref{table:ax:models-2}).


\begin{figure*}% 
\begin{center}
 \begin{tikzpicture}[scale =1,yscale=2, 
    op/.style={draw, fill=white, rounded corners, minimum width=1.25cm, minimum height=0.55cm, 
               inner sep=2pt, font=\small},
    thread/.style={font=\small\bfseries},
    >=stealth
]

% ------------------------------------------------------------
% First group: T_i^1 and T_i^0
% ------------------------------------------------------------

% T_1^1 / T_1^0
\node[thread] (T11) at (0, 0.9) {$T_1^1$};
\node[thread] (T10) at (2.0, 0.9) {$T_1^0$};

\node[op] (s11a) at (0, 0.2) {$\wt(s_1,1)$};
\node[op] (s11b) at (0,-0.5) {$\wt(s_1,2)$};
\node[op] (v11)  at (0,-1.2) {$\wt(v_1,1)$};

\draw[po] (s11a) to (s11b);
\draw[po] (s11b) to (v11);



\node[op] (s10a) at (2.0, 0.2) {$\wt(s_1,1)$};
\node[op] (s10b) at (2.0,-0.5) {$\wt(s_1,2)$};
\node[op] (v10)  at (2.0,-1.2) {$\wt(v_1,0)$};

\draw[po] (s10a) to (s10b);
\draw[po] (s10b) to (v10);

\draw[mo] (s11b) to node[above] {$\mo$} (s10a);
\draw[mo] (s11a) to (s10a);

% T_2^1 / T_2^0
\node[thread] (T21) at (0,-2.4) {$T_2^1$};
\node[thread] (T20) at (2.0,-2.4) {$T_2^0$};

\node[op] (s21a) at (0,-3.1) {$\wt(s_2,1)$};
\node[op] (s21b) at (0,-3.8) {$\wt(s_2,2)$};
\node[op] (v21)  at (0,-4.5) {$\wt(v_2,1)$};

\draw[po] (s21a) to (s21b);
\draw[po] (s21b) to (v21);

\node[op] (s20a) at (2.0,-3.1) {$\wt(s_2,1)$};
\node[op] (s20b) at (2.0,-3.8) {$\wt(s_2,2)$};
\node[op] (v20)  at (2.0,-4.5) {$\wt(v_2,0)$};

\draw[po] (s20a) to (s20b);
\draw[po] (s20b) to (v20);
\draw[mo] (s20b) to node[above] {$\mo$} (s21a);
\draw[mo] (s20a) to (s21a);

% T_3^1 / T_3^0
\node[thread] (T31) at (0,-5.7) {$T_3^1$};
\node[thread] (T30) at (2.0,-5.7) {$T_3^0$};

\node[op] (s31a) at (0,-6.4) {$\wt(s_3,1)$};
\node[op] (s31b) at (0,-7.1) {$\wt(s_3,2)$};
\node[op] (v31)  at (0,-7.8) {$\wt(v_3,1)$};

\draw[po] (s31a) to (s31b);
\draw[po] (s31b) to (v31);

\node[op] (s30a) at (2.0,-6.4) {$\wt(s_3,1)$};
\node[op] (s30b) at (2.0,-7.1) {$\wt(s_3,2)$};
\node[op] (v30)  at (2.0,-7.8) {$\wt(v_3,0)$};

\draw[po] (s30a) to (s30b);
\draw[po] (s30b) to (v30);
\draw[mo] (s30b) to node[above] {$\mo$} (s31a);
\draw[mo] (s30a) to (s31a);
% ------------------------------------------------------------
% Second group: T_xi and T_not xi
% ------------------------------------------------------------

% T_x1 / T_not x1
\node[thread] (Tx1)  at (5.0, 0.9) {$T_{x_1}$};
\node[thread] (Tnx1) at (7.0, 0.9) {$T_{\neg x_1}$};

\node[op] (rv11) at (5.0,0.2) {$\rd(v_1,1)$};
\node[op] (c11)  at (5.0,-0.5) {$\wt(c_1,1)$};
\node[op] (c21)  at (5.0,-1.2) {$\wt(c_2,1)$};

\draw[po] (rv11) to (c11);
\draw[po] (c11) to (c21);

\node[op] (rv10) at (7.0,0.2) {$\rd(v_1,0)$};

% T_x2 / T_not x2
\node[thread] (Tx2)  at (5.0,-2.4) {$T_{x_2}$};
\node[thread] (Tnx2) at (7.0,-2.4) {$T_{\neg x_2}$};

\node[op] (rv21) at (5.0,-3.1) {$\rd(v_2,1)$};
\node[op] (c12)  at (5.0,-3.8) {$\wt(c_1,1)$};

\draw[po] (rv21) to (c12);


\node[op] (rv20) at (7.0,-3.1) {$\rd(v_2,0)$};
\node[op] (c22)  at (7.0,-3.8) {$\wt(c_2,1)$};

\draw[po] (rv20) to (c22);

% T_x3 / T_not x3
\node[thread] (Tx3)  at (5.0,-5.7) {$T_{x_3}$};
\node[thread] (Tnx3) at (7.0,-5.7) {$T_{\neg x_3}$};

\node[op] (rv31) at (5.0,-6.4) {$\rd(v_3,1)$};
\node[op] (c13)  at (5.0,-7.1) {$\wt(c_1,1)$};

\draw[po] (rv31) to (c13);

\node[op] (rv30) at (7.0,-6.4) {$\rd(v_3,0)$};
\node[op] (c23)  at (7.0,-7.1) {$\wt(c_2,1)$};

\draw[po] (rv30) to (c23);
\draw[mo, bend left=45] (c12.west) to (c11);
\draw[mo, bend left=5] (c22.west) to (c21);
\draw[mo, bend left=50] (c13.west) to (c11.west);
\draw[mo, bend right=55] (c23.east) to (c21.east);

% ------------------------------------------------------------
% Final thread T_f
% ------------------------------------------------------------

\node[thread] (Tf) at (10.0,0.9) {$T_f$};

\node[op] (fc1) at (10.0,0.2) {$\rd(c_1,1)$};
\node[op] (fc2) at (10.0,-0.5) {$\rd(c_2,1)$};
\node[op] (fs12) at (10.0,-1.2) {$\rd(s_1,2)$};
\node[op] (fs11) at (10.0,-1.9) {$\rd(s_1,1)$};
\node[op] (fs22) at (10.0,-2.6) {$\rd(s_2,2)$};
\node[op] (fs21) at (10.0,-3.3) {$\rd(s_2,1)$};
\node[op] (fs32) at (10.0,-4.0) {$\rd(s_3,2)$};
\node[op] (fs31) at (10.0,-4.7) {$\rd(s_3,1)$};

%po for Tf
\draw[po] (fc1) to (fc2);
\draw[po] (fc2) to (fs12);
\draw[po] (fs12) to (fs11);
\draw[po] (fs11) to (fs22);
\draw[po] (fs22) to (fs21);
\draw[po] (fs21) to (fs32);
\draw[po] (fs32) to (fs31);

%rf edges

\draw[rf] (c11) to (fc1);
\draw[rf] (c21) to (fc2);
\draw[rf, bend right=8] (s11b) to (fs12);
\draw[rf, bend right=25] (s10a) to (fs11);
\draw[rf, bend left=10] (s20b) to (fs22);
\draw[rf, bend left=10] (s21a) to (fs21);
\draw[rf, bend left=0] (s30b) to (fs32);
\draw[rf, bend left=0] (s31a) to (fs31);
\end{tikzpicture}%
\end{center}
\caption{Execution graph for Table~\ref{tab:example-finegrainedwra}:The trivial $\rf$-edges are omitted}
\label{example-Execution-ra}
\end{figure*}%
 
